\section{Background}

% ============================================================
% 2.1 Agent Lifecycle Management in Multi-Agent Systems
% ============================================================

Modern AI systems are increasingly architected as ecosystems of autonomous agents that perceive, decide, and act without direct human intervention at every step. In such environments, an agent's lifecycle encompasses its creation (or \emph{spawning}), goal-directed execution, potential for self-modification or delegation, and eventual termination. When multiple agents operate concurrently, these lifecycle phases compound: one agent's spawn event becomes another's context boundary, and a delegation handoff becomes a critical audit point. Without explicit lifecycle management, multi-agent systems devolve into untracked action clouds where accountability dissolves \cite{wang2024identity}.

The Agent Lifecycle Protocol (ALP) formalizes six canonical lifecycle events---Genesis, Fork, Migration, Retraining, Succession, and Decommission---each with explicit state-machine semantics \cite{alp2025}. Genesis covers the initial instantiation of an agent with declared identity and mandate. Fork addresses the creation of a child agent whose trust profile inherits partially but non-identically from its parent, enabling controlled capability expansion without wholesale identity copying. Succession manages the orderly transfer of ongoing obligations when an agent retires or is decommissioned. Crucially, ALP embeds a genealogical registry that records both genetic lineage (model, architecture) and epigenetic lineage (configuration, memory, reputation history), supporting queries over an agent's complete ancestry. This registry is the mechanism by which immutable parent chains become operationally tractable: every spawn event is anchored to a traceable prior state.

Dynamic spawning models such as AgentSpawn extend this picture by enabling agents to spawn specialized children at runtime based on observed task complexity \cite{agentspawn2025}. Rather than pre-defining agent topologies, AgentSpawn introduces adaptive spawning policies that trigger child-agent creation when runtime metrics indicate sub-task demand. Crucially, AgentSpawn transfers relevant context and memory to the child agent automatically, preserving epistemic continuity across the spawn boundary. This stands in contrast to static topologies, where agent roles are fixed at design time and cannot adapt to execution-time contingencies.

Accountability in multi-agent organizations has been formally analyzed by \cite{ferrandin2022accountability}, who define an accountability lifecycle comprising account generation, propagation, and corrective action. Their framework maps naturally onto agentic systems: each agent action generates a responsibility account that propagates up the organizational hierarchy until an appropriate corrective authority responds. The key insight is that accountability is not a post-hoc judgment but a structured process that must be engineered into the system's control plane. Their work with the JaCaMo platform demonstrates how accountability mechanisms can be integrated with agent programming abstractions to produce traceable, governable multi-agent organizations.

% ============================================================
% 2.2 Accountability and Provenance in Distributed Systems
% ============================================================

Provenance---the record of an entity's derivation and the causal links between entities---has long been a concern in distributed systems and scientific data management. The W3C PROV standard provides a canonical model for expressing provenance across entities, activities, and agents \cite{moreau2023prov}. In conventional distributed systems, provenance is relatively tractable because computational steps map cleanly onto deterministic software actions. Agentic AI introduces complications: agents operate probabilistically, may invoke external tools, and can exhibit nondeterministic behavior even given identical inputs \cite{openclaw2025forensic}. Recovering a reliable provenance trail therefore requires instrumentation at multiple abstraction layers.

PROV-AGENT extends the W3C PROV model for agentic workflows using the Model Context Protocol (MCP) to capture agent-specific metadata---prompts, responses, intermediate reasoning states, tool invocations---alongside conventional provenance data \cite{provagent2025}. PROV-AGENT addresses the traceability challenge by providing near real-time provenance capture, enabling forensic reconstruction of agent decision chains. The authors demonstrate that their approach detects hallucination propagation across agent chains, where an erroneous intermediate output by one agent propagates undetected to downstream agents. By contrast, systems without provenance instrumentation silently amplify such errors.

Complementing provenance is the question of identity and authorization at scale. The Machine Identity Governance Taxonomy (MIGT) and AI-Identity Risk Taxonomy (AIRT) classify machine identities across six governance domains, emphasizing that non-human identities outnumber human identities in enterprise deployments and require integrated lifecycle management \cite{migt2025}. These taxonomies identify 37 risk subcategories spanning technical, regulatory, and cross-jurisdictional dimensions, underscoring that accountability in agentic systems cannot rely solely on human authentication at spawn time. Rather, each agent must carry a verifiable identity context that persists across delegation hops.

Human-Delegation Provenance (HDP) offers a lightweight, token-based protocol for embedding and verifying human authorization within multi-agent, hierarchical task delegation \cite{hdp2025}. An HDP token ties a human principal, authorization scope, and session to a delegation session; each agent hop appends a cryptographically signed record, forming an append-only provenance chain. Critically, HDP verification is fully offline---it requires only the issuer's Ed25519 public key and the current session ID, with no external registries or trusted third parties. This property makes HDP attractive for agentic deployments where network topology is unreliable or where agents operate across trust boundaries. HDP's chain structure is conceptually equivalent to an immutable parent chain: every downstream action is traceable to its originating human authorization.

The IETF draft on Agent Identity Architecture further codifies a human-anchored identity model in which explicit delegation semantics express the scope, conditions, and authority granted to agents \cite{beyer2025agentidentity}. The architecture separates the durable human identity root from ephemeral agent identities, enabling agents to operate across platforms while remaining attributable to their human principals. This architecture provides the conceptual skeleton on which immutable delegation chains can be implemented: human intent is the immutable anchor, and agent actions are scoped, auditable extensions of that intent.

% ============================================================
% 2.3 Fixed vs. Mutable Delegation Chains
% ============================================================

A delegation chain is the sequence of principal-to-delegate handoffs that connects an originating authority to a terminal action. The literature distinguishes two broad paradigms: \emph{fixed} delegation chains, in which the chain topology is declared at genesis and cannot be altered without explicit re-authorization, and \emph{mutable} delegation chains, in which agents can dynamically re-delegate without obtaining fresh authorization at each hop.

Fixed chains offer maximal auditability: because the chain topology is known in advance, any deviation from the declared path is immediately detectable as a governance violation. The Intelligent AI Delegation framework by \cite{tomasev2026intelligent} argues, however, that rigidity in delegation chains impairs robustness in dynamic environments. Their framework instead advocates \emph{adaptive delegation}: agents continuously monitor capability match and task outcomes, and may re-delegate or escalate when initial assignments fail. This adaptive approach retains auditability by maintaining explicit delegation records but permits chain topology to evolve in response to execution contingencies.

The Recursive Delegation Engine (RDE) operationalizes this tension by modeling delegation as a utility-driven decision problem \cite{rde2025}. An agent evaluates whether to execute a task directly or to delegate it further based on propagated success signals and contextual factors such as execution cost and available context. The resulting delegation topology forms a directed graph in which edges represent permissible delegation relationships. RDE's key audit property is that its decision logic is fully transparent: the utility calculations and multi-armed bandit policies that govern delegation choices are logged and replayable, enabling post-hoc reconstruction of why a particular delegation path was chosen at a particular moment.

From a governance perspective, mutable chains introduce the \emph{delegation drift} problem: as agents re-delegate without re-authorization, the chain of accountability stretches until the originating authority becomes a mere notional anchor rather than a concrete principal. Immutable parent chains address this by anchoring every delegation hop to its immediate parent through a cryptographically signed record. Whether the chain is fixed or mutable at the topology level, the chain of signatures must be append-only: altering a historical record breaks the chain's integrity and is detectable by any downstream verifier.

The trade-off between fixed and mutable chains is fundamentally a trade-off between agility and audit density. Fixed chains minimize audit complexity but cannot adapt to failures or capability gaps. Mutable chains maximize flexibility but require richer instrumentation to maintain provenance. A hybrid approach---in which the chain's topology is mutable but every mutation is recorded immutably---captures both properties. This is the design philosophy adopted by the BX3 Framework.

% ============================================================
% 2.4 Hierarchical Task Decomposition in AI
% ============================================================

Hierarchical task decomposition (HTD) is the process by which a complex goal is recursively broken into sub-goals, each assigned to a specialized agent or skill. HTD is foundational to scalable AI systems because it converts intractable global planning problems into tractable local planning problems: an agent need not plan the entire solution path; it need only decompose the current goal and delegate appropriately.

The Markov Intent Process (MIP) framework by \cite{roijers2020robust} provides a formal foundation for hierarchical planning with policy delegation. MIP agents are composed of specialized skills, each handling a narrow task domain. When a planning goal arrives, the top-level agent analyzes the goal, identifies the skill best suited to the sub-goal, and delegates downward. This creates a dynamic hierarchy in which the delegation chain emerges from the task structure rather than from a fixed design-time assignment. MIP's auditability derives from the hierarchical goal structure itself: because sub-goals are explicitly named and delegated, the full planning tree can be reconstructed from logs.

The Intelligent AI Delegation framework extends HTD beyond pure planning to encompass the full delegation lifecycle \cite{tomasev2026intelligent}. Delegation in this framework is not merely a routing decision but a governance act: it includes the transfer of authority, the declaration of role boundaries, the specification of intent, and the establishment of accountability metrics. This expanded conception of HTD is significant because it treats task decomposition as a provenance event, not merely a computational one. Each decomposition step is recorded, attributable, and subject to verification.

Enterprise-grade agentic systems require coordination among heterogeneous agents (LLM agents, sensor-driven robotics agents, human-in-the-loop oversight agents) operating under formal organizational rules \cite{agenticcommunities2026}. The Agentic Communities framework classifies three tiers of agentic patterns---LLM Agents, Agentic AI, and Agentic Communities---and describes governance structures for the highest tier. In Agentic Communities, inter-agent coordination is mediated by formal roles, protocols, and governance rules that define who may decompose, who may delegate, and what verification obligations each agent carries. This governance layer is where immutable parent chains gain practical significance: without a tamper-evident record of the decomposition hierarchy, there is no reliable basis for audit, correction, or accountability.

The distinction between flat and hierarchical decomposition is important for provenance. In flat decomposition, sub-tasks are assigned independently without explicit parentage relationships; provenance in such systems is typically reconstructed from event logs rather than structural records. In hierarchical decomposition, every sub-task is structurally linked to its parent, forming a tree of accountability that mirrors the task graph. Immutable parent chains make this structural link explicit and tamper-evident.

% ============================================================
% 2.5 The BX3 Framework as Substrate for Immutable Parent Chains
% ============================================================

The BX3 Framework is purpose-built to serve as the operational substrate for immutable parent chains in agentic AI systems. Where prior work in provenance, delegation, and lifecycle management has addressed individual concerns (accountability, identity, hierarchical planning), BX3 integrates these concerns into a unified substrate in which every agent spawn event creates an immutable link to its parent agent, every delegation handoff is a cryptographically signed record, and the complete chain of ancestry is queryable at any point in the system's execution.

BX3's core data model defines an \emph{agent record} that encompasses the agent's identity, its parent reference, its declared mandate, and a signed attestation from its parent confirming the conditions of its creation. This record is append-only: once an agent record is written, it cannot be altered without generating a detectable inconsistency in the chain. The parent reference is not a soft pointer but a cryptographic commitment---the child's record contains a hash of the parent's record at the time of spawn, ensuring that any tampering with the parent's record breaks the child's cryptographic linkage.

The framework's support for recursive spawning builds on the insight that HTD and immutable parent chains are structurally identical: task decomposition creates a parent-child relationship (the decomposer is the parent of the decomposed sub-tasks), and immutable parent chains make this relationship permanent and auditable. In BX3, when an agent decomposes a task and spawns child agents to address sub-components, each child agent carries a verifiable link to its decomposer. The resulting task graph is simultaneously a computational structure (agents solving sub-problems) and a provenance structure (the ancestry of each decision is permanently recorded).

BX3's accountability model draws from the accountability lifecycle framework of \cite{ferrandin2022accountability} but extends it to the recursive spawning case. In BX3, every agent action generates a local accountability record that propagates upward to the agent's parent. The parent aggregates its children's records and adds its own accountability record before propagating further upward. This cascaded aggregation ensures that at any depth in the agent hierarchy, the complete accountability trail from the root agent to the terminal action is reconstructable from the chain of parent links.

The framework's design for fixed vs. mutable chain handling draws from both the ALP genealogical registry and the HDP token structure. BX3 agents are initialized with a declared chain topology: the set of permitted parent references, the maximum delegation depth, and the scope of permissible child-agent mandates. At runtime, the framework records every deviation from the declared topology as a governance event, enabling security and compliance systems to detect unauthorized spawning. This captures the auditability advantages of fixed chains while retaining the flexibility to declare mutable topologies where operational requirements demand it.

By integrating cryptographic parent linking from spawn time onward, the BX3 Framework provides the substrate on which immutable parent chains are not an optional add-on but the default execution environment. Every agent is born into a chain, every delegation extends a chain, and every termination is recorded against the chain. The result is a multi-agent system in which provenance is structural rather than incidental, and accountability is inherited rather than asserted.

% ============================================================
% 2.6 Summary
% ============================================================

The background established in this section reveals a field actively grappling with the governance of autonomous agent systems. Provenance models track agent decisions; lifecycle protocols formalize spawn, succession, and termination; hierarchical decomposition structures task allocation; and cryptographic delegation tokens provide tamper-evident chain integrity. Each thread addresses a specific concern, but their integration into a unified immutable-parent-chain substrate remains an open problem. The BX3 Framework, as described in the subsequent sections, solves this integration by making immutable parent chains the default execution primitive of agentic AI systems, providing structural provenance, inherited accountability, and auditable hierarchical decomposition as first-class system properties.
